% Generated draft: requires mathematical and source review.
\documentclass[11pt]{article}
\usepackage[margin=25mm]{geometry}
\usepackage{fontspec}
\setmainfont{Latin Modern Roman}
\usepackage{amsmath,amssymb,mathtools,unicode-math}
\setmathfont{Latin Modern Math}
\usepackage{enumitem,needspace,xurl,hyperref}
\setlength{\parindent}{0pt}
\setlength{\parskip}{5pt}
\setlength{\emergencystretch}{3em}
\newcommand{\N}{\mathbb{N}}
\newcommand{\Z}{\mathbb{Z}}
\newcommand{\Q}{\mathbb{Q}}
\newcommand{\R}{\mathbb{R}}
\newcommand{\C}{\mathbb{C}}
\newcommand{\F}{\mathbb{F}}
\newcommand{\A}{\mathbb{A}}
\newcommand{\PP}{\mathbb{P}}
\newcommand{\E}{\mathbb{E}}
\newcommand{\eps}{\varepsilon}
\newcommand{\dd}{\,\mathrm{d}}
\begin{document}
\section*{UnsolvedMath ID: 20002285}
\textbf{Source problem code:} catalogue entry
\subsection{Definitions and mathematical statement}
Let \(K\) be a number field and \(\mathcal O_K\) its ring of integers. The source's bound/free-variable collision and final \(\theta/\phi\) mismatch require explicit notation repair. With \(a\) the free variable, the three questions are:
\[
 \begin{array}{ll}
 \text{(a)}&\exists\theta_K(a)\text{ with prefix }\forall t\exists\mathbf y
 \text{ defining }\mathcal O_K\text{ in }K;\\
 \text{(b)}&\text{the same one-scalar-universal prefix defining }\mathbb Z\text{ in }K;\\
 \text{(c)}&\exists\phi_K(a)\text{ with a purely universal prefix defining }\mathbb Z\text{ in }K.
 \end{array}
\]
The universal block in (c) is not stated to contain only one scalar variable. The language, parameters and uniformity across \(K\) are not fully specified; use the intended ring language and retain those choices for review.

The aid reports affirmative answers to all three for \(K=\mathbb Q\), attributed to Koenigsmann. It reports purely universal ten-variable definitions of \(\mathcal O_K\), attributed to Daans, but emphasizes that ten universal variables do not answer the one-universal requirement in (a). Its checked literature did not settle the exact general-number-field questions as stated. Full definitions, references and parameter conventions are absent.

The aid gives a conditional prenex transfer lemma. Suppose \(R\subseteq F\) is defined by
\[
 R(z)\ \Longleftrightarrow\
 \forall\mathbf s\in F^q\;G(z,\mathbf s)\ne0,
\]
and \(A\subseteq R\) is defined inside \(R\) by
\(\forall t\exists\mathbf y\,\Phi(a,t,\mathbf y)\), with \(m\) existential variables and quantifier-free matrix. Standard relativization gives a definition in \(F\) with prefix
\[
 \forall^{q+1}\exists^{q+m}\forall^{qm}.
\]
For clarity an explicit matrix, with \(q\)-tuples \(\mathbf s_0,\mathbf r,\mathbf s_1,\ldots,\mathbf s_m\), is
\[
 \begin{split}
 \forall\mathbf s_0\forall t\exists\mathbf r\exists\mathbf y
 \forall\mathbf s_1\cdots\forall\mathbf s_m\quad
 G(a,\mathbf s_0)\ne0\ \wedge\
 \bigl(&G(t,\mathbf r)=0\ \vee\\
 &(\Phi(a,t,\mathbf y)\wedge
       \bigwedge_{j=1}^mG(y_j,\mathbf s_j)\ne0)\bigr).
 \end{split}
\]
The first conjunct retains membership of the free variable in \(R\); the disjunction handles universal \(t\) outside \(R\). For the reported ten-variable integrality realization this becomes
\(\forall^{11}\exists^{10+m}\forall^{10m}\), provided its matrix has the assumed form. It is not a one-universal formula and therefore does not solve (a) or (b).

Finally the aid reports that polynomial or rational one-parameter maps cannot surject onto \(F^q\), \(q\geq2\), for these infinite fields: their coordinates are algebraically dependent, so their images lie in a proper algebraic locus. This excludes that particular algebraic sweep as a way of compressing universal variables. It says nothing decisive about arbitrary definable arithmetic coding or every logical quantifier reduction.
\subsection{Short English statement}
The source distinguishes one universal variable followed by existentials from an arbitrary universal block. The aid reports the rational-field case, ten-variable integrality progress and an exact relativization bound, while leaving the precise general-number-field requests unresolved. Failure of a one-parameter algebraic sweep does not rule out all definable coding.
\subsection{Sources}
\begin{itemize}
\item Local source record: catalogue entry, catalogue entry; canonical \texttt{aim-logic-notes.json}, \texttt{notes[60]}, Question 1 attributed to Shlapentokh and Eisentr\"ager. Title: ``One-universal definitions of integers in number fields.''
\item AIM notes \emph{Problems related to ``Definability and Decidability Problems in Number Theory''}, workshop 9--13 September 2013, moderated by T. Scanlon, notes by J. Demeyer; supplied reference \url{https://aimath.org/pastworkshops/definabilityinntproblems.pdf}. This draft uses the supplied record without accessing a PDF.
\item Supplied research attributions: Koenigsmann for \(K=\mathbb Q\), Daans for the ten-variable integrality definition. Full bibliographic entries, formula matrices and parameter conventions are not supplied. Attempt 1 gives the explicit relativization and algebraic-dependence obstruction.
\item Registry dated 14 August 2026: ``partially\_solved,'' ``PARTIAL-PROGRESS,'' statement status ``exact''; attempt partial, rigor ``proved,'' novelty confidence low. The machine-generated, non-peer-reviewed aid requires independent verification before citation.
\end{itemize}
\end{document}
